% ============================================================
%  APPENDIX TO CHAPTER 1.  Material NOT in Rudin, imported from
%  Tao I/II, ECNU, USTC, Yu Pin.  Numbered S1.n throughout.
% ============================================================
\rappendix{1}

\begin{roadmap}[How to use this appendix]
Nothing in this appendix is Rudin's text. It is imported from the four
comparison texts and set apart at every level: the running head names the
appendix, section heads carry a reverse-video tag, item numbers are prefixed
\textbf{S}, and every statement names its source.

\medskip
\raggedright
\textbf{Why import anything at all.} Chapter 1 is the one chapter where Rudin
is deliberately incomplete, and he says so. He assumes $\Q$ outright, exiles
the construction of $\R$ to an appendix he tells you to skip, and states his
sharpest result --- that the complete ordered field is unique --- with the
words \emph{``which we will not prove here.''} Each omission is defensible;
each also leaves a question a reader is entitled to have answered.

\medskip
\textbf{What is here.}

\smallskip
\textbf{A. The uniqueness theorem, proved.} Rudin's strongest sentence about
$\R$ is the one he declines to justify. Proving it is what converts his
appendix from ``one construction'' into ``the only possible answer.''

\medskip
\textbf{B. How much is the Archimedean property?} Rudin gets it free from
\lub{} at 1.20(a), so he never has to show it is not automatic. It is not:
there are ordered fields containing infinite elements, and seeing one
explains why Yu Pin must assume Archimedes separately.

\medskip
\textbf{C. Decimals, and a third route to the supremum.} Rudin disposes of
decimals in three lines and says he will never use them. ECNU \emph{defines}
reals as decimals, which makes $0.999\ldots = 1$ a normalisation rule and
turns the supremum into something one computes rather than asserts.

\medskip
\textbf{D. The toolkit Rudin never states.} The real triangle inequality, its
reverse form, the band rule for absolute values, and the
$\varepsilon$-principle --- four facts he uses constantly and states
nowhere.

\medskip
\textbf{E. Below Rudin's starting line.} He assumes $\Q$ in his second
sentence. Tao pays that bill in full, and the itemised cost explains what
Rudin's appendix is an instance of.

\medskip
\textbf{F. $\R$ is not the only completion of $\Q$.} ``Filling the holes''
presumes a way of measuring closeness. Change the absolute value and the
holes move.

\medskip
\textbf{G. Real exponents, worked.} Rudin's Exercises 1.6--1.7 done
properly, including the two steps his hint glosses.

\medskip
\textbf{How to read it.} Section A belongs with Rudin's own appendix, and is
best read straight after it. Section B can be read as soon as 1.20 is done,
C alongside 1.22, D before Chapter 3, and G with Exercises 1.6--1.7. E and F
are optional and answer ``what did the assumptions cost?'' The exercises are
the last section.
\end{roadmap}

% ============================================================
\ssection{Supplement A --- The uniqueness theorem, proved}

\begin{roadmap}[Why this supplement belongs to Chapter 1]
\raggedright
\textbf{What Rudin does.} Theorem 1.19 asserts that an ordered field $\R$
with the least-upper-bound property exists and contains $\Q$ as a subfield;
his appendix builds one from Dedekind cuts. At the end of that appendix he
writes: \emph{``It is a fact, which we will not prove here, that any two
ordered fields with the least-upper-bound property are isomorphic''} --- and
concludes that 1.19 \emph{characterizes} $\R$ completely.

\smallskip
\textbf{The gap.} That sentence is the load-bearing one for the entire book,
and it arrives unsupported. Everything after Chapter 1 uses only the
\emph{properties} of $\R$ and never the cuts; the licence for that is
precisely uniqueness. Without it, a sceptic could ask whether the theorems of
Chapters 2--11 are theorems about \emph{Dedekind's} real numbers rather than
about the real numbers.

\smallskip
\textbf{Where the gap bites.} Immediately and permanently. It is why Rudin
can tell the reader, in his Preface, to skip the appendix; why Tao may build
$\R$ from Cauchy sequences and get the same object; why ECNU may define reals
as infinite decimals; and why the six equivalent completeness axioms of the
Chapter 1 box are equivalent \emph{as descriptions of one thing} rather than
as six different subjects.

\smallskip
\textbf{What this section supplies.} The proof. It is not hard --- it is the
density of $\Q$ used four times --- but it is the argument that retires the
construction for good.
\end{roadmap}

\sitem{S1.1}{Theorem}
\emph{Let $F_1$ and $F_2$ be ordered fields with the least-upper-bound
property. Then there is exactly one order-preserving field isomorphism
$\varphi : F_1 \to F_2$, and it fixes the rationals.}
\provenance{Nobody's. Rudin states it and declines to prove it; Tao never
introduces the notion of an ordered field at all; Yu Pin defers it (p.\,35,
``as for uniqueness, we will discuss it in passing later'') and delivers only
the metric version --- the completion of a metric space is unique up to
isometry (p.\,124). The proof below is supplied here.}

\begin{proof}
Each of $F_1$, $F_2$ contains an isomorphic copy of $\Q$, obtained by sending
$n \mapsto 1 + \cdots + 1$ and then $m/n \mapsto m \cdot n^{-1}$; the field
axioms make this injective and order-preserving, and we identify $\Q$ with
its image in each. By 1.20, applied in each field, both are archimedean and
$\Q$ is dense in each.

For $x \in F_1$ put
\[ Q_x = \{\,q \in \Q : q < x\,\}, \qquad
   \varphi(x) = \sup\nolimits_{F_2} Q_x . \]
The set $Q_x$ is nonempty and bounded above in $F_2$: by the archimedean
property there are integers $m < x < n$, so $m \in Q_x$ and $n$ bounds it.
Hence the supremum exists in $F_2$ by hypothesis, and $\varphi$ is defined.

\emph{$\varphi$ fixes $\Q$ and preserves order.} If $r \in \Q$ then
$Q_r = \{q \in \Q : q < r\}$ has supremum $r$ in $F_2$, so $\varphi(r)=r$. If
$x < y$ in $F_1$, density gives rationals with $x < q < r < y$; then every
element of $Q_x$ is below $q$ while $r \in Q_y$, so
$\varphi(x) \le q < r \le \varphi(y)$. In particular $\varphi$ is injective.

\emph{$\varphi$ is additive.} $Q_x + Q_y \subset Q_{x+y}$ gives
$\varphi(x)+\varphi(y) \le \varphi(x+y)$. Conversely, if $q < x+y$ is
rational, density supplies rationals $q_1 < x$ and $q_2 < y$ with
$q < q_1+q_2$ --- choose $q_1$ within $\tfrac12(x+y-q)$ of $x$ and $q_2$
likewise --- so $q \le \varphi(x)+\varphi(y)$. Taking the supremum over such
$q$ gives the reverse inequality.

\emph{$\varphi$ is multiplicative.} For $x,y>0$ the same two-sided argument
applies to products of positive rationals; the general case follows from
additivity, since $\varphi(-x) = -\varphi(x)$.

\emph{$\varphi$ is onto.} Given $z \in F_2$, set
$x = \sup_{F_1}\{q \in \Q : q < z\}$, which exists for the same reason as
before. Then $Q_x$ and $\{q \in \Q : q < z\}$ have the same elements, by
density in each field, so $\varphi(x) = z$.

\emph{Uniqueness of $\varphi$.} Any order-preserving field isomorphism fixes
$1$, hence fixes $\Q$; and an order-preserving map fixing a dense set is
determined, since $\psi(x)$ must be an upper bound of $\psi(Q_x) = Q_x$ and
below every rational above $x$.
\end{proof}

\begin{proofwalk}[How the proof flows]
  \item \textbf{Choose the method.} An isomorphism must be \emph{built}, and
        the only structure both fields are known to share is $\Q$ and the
        supremum operation. So the map must be defined by naming each $x$
        with rationals and reassembling it on the other side: pattern 3 of
        Chapter 1's strategy index, transport to the supremum.
  \item \textbf{Buy the supremum.} $\varphi(x) = \sup Q_x$ needs $Q_x$
        nonempty and bounded --- exactly the two tickets 1.10 demands, and
        exactly what the archimedean property (1.20(a)) supplies. Note this
        is the first use of \lub{}, and it is used to \emph{define} the map,
        not to prove anything about it.
  \item \textbf{Get order for free.} Density (1.20(b)) separates any two
        points by two rationals, and two rationals are enough to separate
        their images. Injectivity comes with it.
  \item \textbf{Do the algebra by squeezing.} Additivity and multiplicativity
        are each proved by two inequalities: one inclusion of sets gives
        ``$\le$,'' and density gives ``$\ge$.'' There is no computation
        anywhere --- only the fact that rationals approximate from below as
        closely as one likes.
  \item \textbf{Run the construction backwards for surjectivity.} The same
        formula with the roles of $F_1$ and $F_2$ exchanged produces the
        preimage. Symmetry of the hypotheses is what makes this free.
  \item \textbf{Audit.} The least-upper-bound property is spent twice, once
        in each field, both times to define a map. The archimedean property
        and density of $\Q$ are spent repeatedly --- but both are
        \emph{consequences} of \lub{} by 1.20, so the hypothesis list is
        exactly Rudin's. Nothing about cuts, decimals or Cauchy sequences
        appears anywhere, which is the whole point.
\end{proofwalk}

\begin{deeper}[What the theorem licenses]
Read the theorem as a permission slip. It says the phrase ``the real
numbers'' denotes a single object, so that:

\smallskip
\emph{Constructions are disposable.} Rudin's cuts, Tao's Cauchy sequences,
and ECNU's infinite decimals are three routes to the same destination.
Rudin's advice in the Preface --- that beginning with a construction is
``pedagogically unsound (though logically correct)'' --- is only defensible
because of this theorem: you may skip the construction precisely because the
construction does not matter.

\smallskip
\emph{Axiom choices are disposable too.} Chapter 1's box compared four
axiomatisations --- Rudin's \lub{}, Yu Pin's nested intervals, Tao's Cauchy
completeness, ECNU's six-cycle. Uniqueness is what makes those four
descriptions of one field rather than four fields.

\smallskip
\emph{And it is sharp.} Drop any hypothesis and the conclusion fails.
Without \lub{} there are many non-isomorphic archimedean ordered fields ---
$\Q$ itself, and every subfield of $\R$. Without the ordering there are
field isomorphisms of $\C$ that are wildly discontinuous. The next section
drops the archimedean property, which \lub{} was silently supplying, and
shows what escapes.
\end{deeper}

% ============================================================
\ssection{Supplement B --- How much is the Archimedean property?}

\begin{roadmap}[Why this supplement belongs to Chapter 1]
\raggedright
\textbf{What Rudin does.} Theorem 1.20(a) derives the archimedean property
--- for $x>0$ and any $y$ there is a positive integer $n$ with $nx>y$ ---
from the least-upper-bound property, in four lines. His margin remark notes
that this is a theorem, not an axiom.

\smallskip
\textbf{The gap.} Deriving a property is not the same as showing it has
content. Because Rudin obtains Archimedes free of charge, a reader has no
way to see that it is a genuine restriction --- that ordered fields exist in
which some element exceeds every integer, and that such fields therefore
cannot have the \lub{} property.

\smallskip
\textbf{Where the gap bites.} At Chapter 1's own comparison box, where Yu Pin
axiomatises $\R$ as field $+$ order $+$ \emph{archimedean} $+$ nested
intervals, and the natural question --- why does he need the third
hypothesis when Rudin does not? --- has no answer inside Rudin. It bites
again at S2.7 in the Chapter 2 appendix, where recovering $\sup$ from Cauchy
completeness needs Archimedes as a separate hypothesis, for exactly the
reason exhibited here.

\smallskip
\textbf{What this section supplies.} An explicit non-archimedean ordered
field, small enough to hold in the hand, together with the proof that no such
field can have the least-upper-bound property.
\end{roadmap}

\sitem{S1.2}{Example}
\emph{(An ordered field with infinite elements.) Let $\R(t)$ be the field of
rational functions $p(t)/q(t)$ with real coefficients, and declare
$f > 0$ when $f(x) > 0$ for all sufficiently large real $x$ --- equivalently,
when the leading coefficients of $p$ and $q$ have the same sign. This makes
$\R(t)$ an ordered field in which}
\[ t > n \qquad \text{for every positive integer } n . \]
\provenance{Supplied here: an explicit non-archimedean ordered field is
absent from all four comparison texts. Yu Pin's axiom list (field, order,
archimedean, nested intervals) is built to expose the possibility --- he
notes at p.\,26 that (F), (O) and (A) alone are satisfied by $\Q$ --- but he
gives no example in which (A) fails.}

\begin{unpack}[Checking the order, and reading the element $t$]
That the relation is an order needs only that a nonzero rational function has
constant sign near $+\infty$ --- it has finitely many zeros --- and that the
sum and product of two eventually-positive functions are eventually positive.
Both are immediate, and they are the two clauses of Definition 1.17.

Now look at $t$. For any integer $n$ the function $t - n$ is eventually
positive, so $t > n$: the single element $t$ exceeds every integer. Its
reciprocal $1/t$ is then a positive element smaller than every $1/n$ --- a
genuine infinitesimal. Neither is a limit of anything; they are simply
elements of the field, and the field's arithmetic is entirely ordinary.

The moral is that ``$\R$ has no infinite or infinitesimal elements'' is not a
consequence of being an ordered field. It has to be imposed, and Rudin
imposes it without naming it, by assuming \lub{}.
\end{unpack}

\sitem{S1.3}{Theorem}
\emph{An ordered field with the least-upper-bound property is archimedean.
Consequently $\R(t)$ does not have the least-upper-bound property: the set of
positive integers is bounded above in $\R(t)$ but has no supremum.}

\begin{proof}
This is Rudin's 1.20(a) restated, and his proof is the argument. Suppose the
set $\N$ of positive integers were bounded above in an ordered field with the
\lub{} property, and let $\alpha = \sup \N$. Since $\alpha - 1 < \alpha$, the
element $\alpha-1$ is not an upper bound, so $n > \alpha-1$ for some
$n \in \N$; but then $n+1 > \alpha$ and $n+1 \in \N$, contradicting that
$\alpha$ is an upper bound.

Hence in such a field $\N$ is unbounded, which is the archimedean property.
In $\R(t)$ the element $t$ bounds $\N$ above, so $\R(t)$ cannot have the
\lub{} property.
\end{proof}

\begin{deeper}[Three fields, and what separates them]
Line up the examples Chapter 1 now has:

\smallskip
\begin{tabular}{@{}l c c p{0.44\linewidth}@{}}
\toprule
\textbf{field} & \textbf{arch.} & \textbf{\lub{}} & \textbf{what it lacks} \\
\midrule
$\Q$      & yes & no  & the sup of $\{p : p^2<2\}$ (1.1) \\
$\R(t)$   & no  & no  & any bound on $\N$ \\
$\R$      & yes & yes & --- \\
\bottomrule
\end{tabular}

\smallskip
$\Q$ and $\R(t)$ fail to be $\R$ for genuinely different reasons, and the
difference is worth holding on to. $\Q$ has the right \emph{scale} --- it is
archimedean, its integers are unbounded --- and is merely full of holes.
$\R(t)$ has no holes of $\Q$'s kind at all in the relevant sense; it fails
earlier and more crudely, by containing elements no integer can reach.

This is why Yu Pin's axiom list needs four items where Rudin's needs one
extra property. Nested intervals alone does not imply Archimedes --- one can
have a non-archimedean field in which every nested sequence of intervals has
a common point --- so he must assume it. Rudin's \lub{} implies both at once,
which is economical but conceals that two separate things were bought.
\end{deeper}

% ============================================================
\ssection{Supplement C --- Decimals, and a third route to the supremum}

\begin{roadmap}[Why this supplement belongs to Chapter 1]
\raggedright
\textbf{What Rudin does.} Theorem 1.22 disposes of decimals in three lines:
given $x>0$, take the largest integer $n_0 \le x$, then the largest $n_1$
with $n_0 + n_1/10 \le x$, and so on; the resulting decimal converges to $x$.
He then writes that \emph{``we shall never use decimals''} and does not
mention them again.

\smallskip
\textbf{The gap.} Two, in opposite directions. Forwards: his appendix builds
$\R$ from cuts and stops, never reconnecting the constructed field to the
notation every reader has used since childhood --- so ``what \emph{is}
$0.999\ldots$?'' is never answered. Backwards: he presents \lub{} as an
axiom and its construction as a one-off, leaving the impression that a
supremum is an object one may assert but not compute.

\smallskip
\textbf{Where the gap bites.} At 1.22 itself, where the reader is told the
decimal expansion exists and is unique but not what makes the uniqueness
work; at the end of the appendix, where cuts are abandoned without ever
being converted into anything familiar; and in every later argument that
approximates a real by rationals --- 1.20(b), 1.21, 3.20 --- each of which
rebuilds an approximation scheme that 1.22 could have supplied once.

\smallskip
\textbf{What this section supplies.} ECNU's development, which runs the other
way: reals \emph{are} infinite decimals, the awkward identification
$0.999\ldots = 1$ is a normalisation rule chosen in advance, and the
supremum is produced by an algorithm --- one digit at a time --- rather than
asserted.
\end{roadmap}

\sitem{S1.4}{Definition}
\emph{(Reals as decimals.) A real number is an infinite decimal
$a_0.a_1a_2a_3\cdots$; it is rational when the expansion terminates or
eventually repeats and irrational otherwise. To make the representation
unique, every terminating expansion is rewritten with a tail of nines:
for $a_n \ne 0$,}
\[ a_0.a_1\cdots a_n \;=\; a_0.a_1\cdots(a_n-1)999\cdots , \]
\emph{and an integer $a_0$ is written $(a_0-1).999\cdots$. Order is then
lexicographic: $x > y$ if $a_0 > b_0$, or if the digits agree up to place
$l$ and $a_{l+1} > b_{l+1}$.}
\provenance{ECNU \cjk{定义}{Def} 1 (bk.\ pp.\,1--2), where this is the
definition of $\R$, not a theorem about it.}

\begin{pitfall}[$0.999\ldots = 1$ is a convention, not a paradox]
Students meet this identity as a puzzle. In ECNU's development it is not a
puzzle at all: $0.999\ldots$ is the \emph{official name} of $1$, chosen so
that the lexicographic order of Definition S1.4 is a genuine total order.

Why the choice matters: without a convention, $2.001$ and $2.000999\ldots$
are two names for one number, and lexicographic comparison of names would
declare the first larger. Normalising every terminating decimal to its
nine-tail leaves exactly one name per number, and comparison of names then
computes comparison of numbers. Note the convention runs \emph{opposite} to
intuition --- the canonical form is the infinite one.

Rudin sidesteps all of this by refusing to use decimals. His 1.22 asserts
existence and uniqueness of the expansion and moves on; the uniqueness there
is bought by taking the \emph{largest} $n_k$ at each stage, which produces
the terminating form and so is the opposite normalisation. Neither choice is
better --- but a reader who has seen only one will find the other's
$0.999\ldots$ baffling.
\end{pitfall}

\sitem{S1.5}{Definition}
\emph{(Deficient and excess approximations.) For $x = a_0.a_1a_2\cdots \ge 0$
put}
\[ x_n = a_0.a_1\cdots a_n, \qquad \bar x_n = x_n + 10^{-n} , \]
\emph{so that $x_n \le x < \bar x_n$ with $\bar x_n - x_n = 10^{-n}$. The
$x_n$ increase and the $\bar x_n$ decrease.}
\provenance{ECNU \cjk{定义}{Def} 2 (bk.\ p.\,2).}

\sitem{S1.6}{Theorem}
\emph{$x > y$ if and only if $x_n > \bar y_n$ for some $n$. Consequently, if
$x < y$ there is a rational $r$ with $x < r < y$.}
\provenance{ECNU \cjk{命题}{Prop} and \cjk{例}{Ex} 1 (bk.\ pp.\,2--3).}

\begin{proof}
If $x_n > \bar y_n$ then $x \ge x_n > \bar y_n > y$. Conversely, if $x>y$
then $x - y > 0$, and since $\bar x_n - x_n = 10^{-n} \to 0$ the sandwich
$x_n \le x$, $y < \bar y_n$ forces $x_n > \bar y_n$ once $10^{-n}$ is smaller
than the gap.

For density, apply this with the roles reversed: $x<y$ gives an $n$ with
$\bar x_n < y_n$, and
\[ r = \tfrac12\bigl(\bar x_n + y_n\bigr) \]
is rational --- being the average of two terminating decimals --- with
$x \le \bar x_n < r < y_n \le y$.
\end{proof}

\begin{deeper}[Two proofs of density, and what each costs]
Compare with Rudin's 1.20(b). He proves density from the archimedean
property: stretch the interval by an integer $n$ until $n(y-x)>1$, catch an
integer $m$ in the stretched gap, and divide back. The catching step needs
that a set of integers bounded below has a least element --- well-ordering
--- which he uses without comment.

The proof above needs none of that. It truncates both numbers at a common
depth and averages, and the only fact used is that the truncations sandwich
the number within $10^{-n}$. It is also \emph{constructive}: it names the
rational, whereas Rudin's argument asserts that a suitable $m$ exists.

The trade is that this proof is available only once reals are decimals.
Rudin's works in any archimedean ordered field, which is what he needs, since
at 1.20 he has no decimals and does not want any.
\end{deeper}

\sitem{S1.7}{Theorem}
\emph{(The supremum, computed.) Every nonempty set $S \subset \R$ that is
bounded above has a least upper bound, and its decimal digits can be
generated one at a time.}
\provenance{ECNU \cjk{确界原理}{sup principle} \cjk{定理}{Thm} 1.1
(bk.\ pp.\,5--7), proved from the decimal definition.}

\begin{proof}
Assume $S$ contains a nonnegative number. Choose a nonnegative integer $n$
such that $x < n+1$ for every $x \in S$ while some $a_0 \in S$ has
$a_0 \ge n$; this is possible because $S$ is bounded above and nonempty.

Divide $[n, n+1)$ into ten equal parts. Among the ten, there is a largest
digit $n_1$ for which some element of $S$ reaches $n.n_1$, and then
\[ x < n.n_1 + 10^{-1} \quad\text{for all } x \in S,
   \qquad a_1 \ge n.n_1 \quad\text{for some } a_1 \in S. \]
Divide $[\,n.n_1,\ n.n_1 + 10^{-1})$ into ten parts and repeat. Inductively
one obtains digits $n_1, n_2, \dots$ with
\begin{equation}
  x < n.n_1\cdots n_k + 10^{-k}\ \ (x \in S), \qquad
  a_k \ge n.n_1\cdots n_k \ \ \text{for some } a_k \in S. \tag{S1.1}
\end{equation}
Let $\eta = n.n_1n_2n_3\cdots$ be the real number those digits name.

$\eta$ is an upper bound: if some $x \in S$ had $x > \eta$, then by S1.6
$x_k > \bar\eta_k = n.n_1\cdots n_k + 10^{-k}$ for some $k$, contradicting the
first half of (S1.1).

$\eta$ is least: if $\alpha < \eta$, then by S1.6 there is $k$ with
$\eta_k > \bar\alpha_k$, and the second half of (S1.1) supplies $a_k \in S$
with $a_k \ge \eta_k > \bar\alpha_k \ge \alpha$, so $\alpha$ is not an upper
bound.
\end{proof}

\begin{proofwalk}[How the proof flows]
  \item \textbf{Choose the method.} Unlike Rudin, we may not assume the
        supremum exists --- it is the conclusion. So it must be
        \emph{constructed}, and the only construction material available is
        decimal digits. The proof is therefore an algorithm: decide the
        digits one at a time.
  \item \textbf{Fix the integer part.} Boundedness above gives an integer
        ceiling; nonemptiness gives a floor. Both hypotheses of 1.10 are
        spent here, in the first sentence, exactly as in Rudin's 1.21.
  \item \textbf{Choose each digit greedily.} At each stage take the largest
        digit that some element of $S$ still reaches. ``Largest that is still
        attained'' is what makes the result least; ``still an upper bound''
        is what makes it an upper bound. The two clauses of (S1.1) are the
        two halves of Definition 1.8, maintained as loop invariants.
  \item \textbf{Let the algorithm run forever.} The digits define a real
        number by S1.4 --- and this is the step that would fail over $\Q$,
        since the digit string need not repeat.
  \item \textbf{Verify both clauses.} Each is checked by S1.6, which converts
        a comparison of reals into a comparison of finite truncations. The
        invariants were designed so that this conversion closes the argument
        immediately.
  \item \textbf{Audit.} Boundedness and nonemptiness are spent once each, in
        step 2. The decimal definition of $\R$ is spent once, in step 4.
        Nothing else enters --- in particular no cuts, no Cauchy sequences,
        and no appeal to \lub{}, which is what makes this a genuine
        \emph{third} construction rather than a restatement.
\end{proofwalk}

\begin{center}
\begin{tikzpicture}[x=1mm, y=1mm]
  \draw[axis] (0,26) -- (104,26);
  \foreach \x in {0,10,...,100}{\draw[thin] (\x,24.6) -- (\x,27.4);}
  \node[mlbl, anchor=north] at (0,23) {$n$};
  \node[mlbl, anchor=north] at (100,23) {$n{+}1$};
  \fill[black!14] (60,25) rectangle (70,27);
  \node[mlbl, anchor=south] at (65,28) {$n_1=6$};
  % stage 2
  \draw[axis] (0,10) -- (104,10);
  \foreach \x in {0,10,...,100}{\draw[thin] (\x,8.6) -- (\x,11.4);}
  \fill[black!14] (30,9) rectangle (40,11);
  \node[mlbl, anchor=south] at (35,12) {$n_2=3$};
  \node[mlbl, anchor=north] at (0,7) {$n.6$};
  \node[mlbl, anchor=north] at (100,7) {$n.7$};
  \draw[guide, {Stealth[length=3pt]}-] (60,23.4) -- (60,13);
  \draw[guide, {Stealth[length=3pt]}-] (70,23.4) -- (100,13);
  \node[lbl, anchor=west] at (0,-2)
    {each stage magnifies the chosen tenth and picks the next digit};
\end{tikzpicture}
\end{center}
\figcap{The supremum computed digit by digit: at each stage the interval is
cut into ten, and the algorithm takes the rightmost part that some element of
$S$ still reaches. The two conditions (S1.1) --- still an upper bound, still
attained --- are the two clauses of Definition 1.8 carried along as
invariants, and the digits they force out spell $\sup S$.}

% ============================================================
\ssection{Supplement D --- The toolkit Rudin never states}

\begin{roadmap}[Why this supplement belongs to Chapter 1]
\raggedright
\textbf{What Rudin does.} He proves the triangle inequality twice --- at
1.33(e) for complex numbers and at 1.37(e) for $\R^k$ --- and defines
neighbourhoods once, at 2.18(a), as open balls.

\smallskip
\textbf{The gap.} The real-line versions of both are missing. The real
triangle inequality never appears as a standalone statement; the reverse
form $\bigl||a|-|b|\bigr| \le |a-b|$ appears nowhere at all; the rule that
converts $|a| < h$ into a two-sided band is never written down; and the
$\varepsilon$-principle --- if $a < b + \varepsilon$ for every
$\varepsilon>0$ then $a \le b$ --- is used in almost every proof from
Chapter 1 onward without ever being stated.

\smallskip
\textbf{Where the gap bites.} Every $\varepsilon$-argument in the book opens
by unpacking an absolute value into a band and closes by letting
$\varepsilon$ shrink. Rudin performs both moves silently --- at 1.20, 1.21,
3.2, 3.3, 4.2 --- so a reader learns them by imitation rather than by
statement, and the reverse triangle inequality, which is what one needs to
bound a difference of norms from below, has to be rediscovered each time it
is wanted.

\smallskip
\textbf{What this section supplies.} The four facts, stated once, with the
proof of the triangle inequality that derives it from the band rule rather
than by cases.
\end{roadmap}

\sitem{S1.8}{Theorem}
\emph{(Absolute value.) For real $a,b$ and $h>0$:}
\begin{enumerate}[label=(\alph*), leftmargin=2.2em, itemsep=2pt, topsep=3pt]
  \item $-|a| \le a \le |a|$;
  \item $|a| < h$ \emph{if and only if} $-h < a < h$;
  \item $\bigl|\,|a|-|b|\,\bigr| \;\le\; |a \pm b| \;\le\; |a|+|b|$.
\end{enumerate}
\provenance{ECNU \S1.2, six properties of $|\cdot|$ (bk.\ p.\,3).}

\begin{proof}
(a) and (b) are the definition read in the two cases $a \ge 0$ and $a<0$.

For (c), add the two instances of (a):
\[ -\bigl(|a|+|b|\bigr) \;\le\; a+b \;\le\; |a|+|b| , \]
and apply (b) in the reverse direction to read that band as
$|a+b| \le |a|+|b|$. Replacing $b$ by $-b$ gives the same bound for
$|a-b|$. For the left-hand inequality, write $|a| = |(a-b)+b| \le |a-b|+|b|$,
so $|a|-|b| \le |a-b|$; exchanging $a$ and $b$ gives
$|b|-|a| \le |a-b|$, and the two together are the claim.
\end{proof}

\begin{unpack}[Why (b) is the load-bearing one]
Parts (a) and (c) are the famous ones, but (b) is the one used most often and
the one Rudin never states. It is the rule that turns a statement about a
\emph{distance} into a statement about \emph{two inequalities}, and back
again --- and the proof above shows the ``back again'' direction doing real
work: the triangle inequality is proved by producing a band and then reading
it as an absolute value, with no case analysis anywhere.

Every $\varepsilon$-$\delta$ argument in the book is an application. When
Rudin writes ``$|s_n - s| < \varepsilon$'' and then concludes
``$s - \varepsilon < s_n < s + \varepsilon$,'' that is (b), used without
comment.
\end{unpack}

\sitem{S1.9}{Theorem}
\emph{(The $\varepsilon$-principle.) If $a < b + \varepsilon$ for every
$\varepsilon > 0$, then $a \le b$. Equivalently, if
$|a - b| < \varepsilon$ for every $\varepsilon>0$, then $a = b$.}
\provenance{ECNU \cjk{例}{Ex} 2 (bk.\ p.\,3) and \cjk{习题}{Ex} 1.1 \#3.}

\begin{proof}
Suppose $a > b$. Then $\varepsilon = a - b$ is positive, and the hypothesis
applied to this particular $\varepsilon$ gives $a < b + (a-b) = a$, which is
absurd.
\end{proof}

\noindent
The proof is three lines and the statement is used everywhere, which is
exactly why it deserves a name.\kn{Watch for it in Rudin. The last step of
1.11's proof, of 1.21, of 3.2(b), and of every uniqueness argument in
Chapter 3 is this principle: having shown a quantity is smaller than every
positive number, conclude it is zero. Rudin never separates the move from
its instances, so it reads as a trick repeated rather than a theorem
applied.}

\sitem{S1.10}{Definition}
\emph{(Neighbourhood system.) For $a \in \R$ and $\delta>0$ write}
\[ U(a;\delta) = \{x : |x-a|<\delta\}, \qquad
   U^{\circ}(a;\delta) = \{x : 0<|x-a|<\delta\}, \]
\emph{and $U_{+}(a;\delta) = [a, a+\delta)$,
$U_{-}(a;\delta) = (a-\delta, a]$ for the one-sided versions;
$U(+\infty) = \{x : x>M\}$, $U(-\infty) = \{x : x<-M\}$ and
$U(\infty) = \{x : |x|>M\}$ for $M$ large.}
\provenance{ECNU \S2.1 (bk.\ pp.\,4--5).}

\begin{deeper}[One notation, and the four definitions it collapses]
Rudin introduces neighbourhoods once, at 2.18(a), and only the two-sided
ball. Every limit notion he later defines is then written out separately:
4.1 for $\lim_{x\to p}$, 4.25 for one-sided limits, 4.33 for limits at
$\pm\infty$. The three definitions look unrelated on the page.

With the system above they are one sentence. In every case,
$\lim f = A$ means
\[ \text{for each } U(A;\varepsilon) \text{ there is }
   U^{\circ}(a;\delta) \text{ with } f\bigl(U^{\circ}(a;\delta)\bigr)
   \subset U(A;\varepsilon), \]
and the four notions differ only in which neighbourhood of $a$ is used:
punctured, right-hand, left-hand, or a neighbourhood of $+\infty$. The
Chapter 4 appendix reaches the same unification from the other direction,
through Tao's ``limit along a subset'' (S4.2); ECNU gets there by naming the
neighbourhoods instead of the index sets.
\end{deeper}

% ============================================================
\ssection{Supplement E --- Below Rudin's starting line}

\begin{roadmap}[Why this supplement belongs to Chapter 1]
\raggedright
\textbf{What Rudin does.} The second sentence of the book fixes his starting
line: \emph{``we shall not, however, enter into any discussion of the axioms
that govern the arithmetic of the integers, but assume familiarity with the
rational numbers.''} Everything in Chapter 1 is built on $\Q$, taken as given.

\smallskip
\textbf{The gap.} A reader cannot tell whether that assumption is cheap or
expensive. Is $\Q$ a few lines of set theory, or a serious construction? And
if it is serious, is Rudin's Dedekind-cut appendix a special trick or an
instance of something standard?

\smallskip
\textbf{Where the gap bites.} At the appendix, which looks like a
one-off contrivance --- why \emph{sets of rationals}, of all things? --- and
at every use of induction in the book, since Rudin never says what property
of $\N$ makes induction legitimate.

\smallskip
\textbf{What this section supplies.} Tao's construction of $\N$, $\Z$ and
$\Q$, priced honestly; and the observation that his three constructions and
Rudin's cuts are the same manoeuvre performed four times.
\end{roadmap}

\sitem{S1.11}{Definition}
\emph{(The Peano axioms.) $\N$ is a set with a distinguished element $0$ and
a successor operation $n \mapsto n{+}{+}$ satisfying: $0$ is not a successor;
the successor operation is injective; and if $P(0)$ holds and $P(n)$ implies
$P(n{+}{+})$ for every $n$, then $P(n)$ holds for all $n$.}
\provenance{Tao I, Axioms 2.1--2.5 (bk.\ pp.\,16--21).}

\begin{deeper}[What Rudin's one sentence costs]
Tao pays the bill in full, and the itemisation is instructive.

\smallskip
\emph{Even addition is a theorem first.} Before $+$ can be \emph{defined} by
recursion one must prove that recursive definitions are legitimate at all ---
Tao's Proposition 2.1.16. Induction is what licenses it. This is the answer
to ``what makes $\N$ special'': not that its elements are small, but that it
supports recursion.

\smallskip
\emph{$\Z$ is a quotient.} An integer is a \emph{formal difference}
$a \,\text{---}\, b$ of naturals, with $a\,\text{---}\,b = c\,\text{---}\,d$
declared to mean $a+d = c+b$. One then checks that $+$, $\times$ and $<$ are
well defined on the classes.

\smallskip
\emph{$\Q$ is another quotient.} A rational is a \emph{formal quotient}
$a /\!/ b$ with $b \ne 0$, and $a/\!/b = c/\!/d$ means $ad = cb$. Same
checking again.

\smallskip
So Rudin's ``assume familiarity with the rational numbers'' is roughly
thirty-five pages, two quotient constructions and a theorem about recursion.
His Preface defends skipping it as avoiding what is ``pedagogically unsound
(though logically correct),'' and after S1.1 --- uniqueness --- that defence
is airtight: the construction cannot matter, so it may be skipped.
\end{deeper}

\begin{unpack}[One manoeuvre, performed four times]
Look at what $\Z$, $\Q$, $\R$ have in common as constructions:

\smallskip
\begin{tabular}{@{}l l l@{}}
\toprule
\textbf{to build} & \textbf{formal symbol} & \textbf{identified when} \\
\midrule
$\Z$ from $\N$ & $a \,\text{---}\, b$   & $a+d = c+b$ \\
$\Q$ from $\Z$ & $a /\!/ b$             & $ad = cb$ \\
$\R$ from $\Q$ (Tao) & $\operatorname{LIM} a_n$ & $\lim|a_n-b_n| = 0$ \\
$\R$ from $\Q$ (Rudin) & a cut $\alpha$ & --- (the set \emph{is} the number) \\
\bottomrule
\end{tabular}

\smallskip
The first three are identical in shape: introduce a symbol for the operation
you wish you could perform, declare when two symbols name the same thing,
and prove the arithmetic descends to the classes. Tao says so explicitly ---
$\operatorname{LIM}$ is introduced as ``a new formal symbol, similar to
$\text{---}$ and $/\!/$'' --- and adds that once formal limits are shown to
be actual limits, they are never needed again.

Rudin's cuts are the odd one out: a cut is not a symbol standing for
something absent but a set that \emph{is} the number. That is why his
appendix needs no equivalence classes and why Step 1 spends so long on
property (III) --- the normalisation that makes the representation unique is
built into the definition rather than quotiented out afterwards.
\end{unpack}

% ============================================================
\ssection{Supplement F --- $\R$ is not the only completion of $\Q$}

\begin{roadmap}[Why this supplement belongs to Chapter 1]
\raggedright
\textbf{What Rudin does.} Section 1.1 exhibits the gap in $\Q$ --- the set
$\{p : p^2<2\}$ with no rational supremum --- and 1.19 supplies $\R$ as the
field that fills it. The narrative is: $\Q$ has holes, $\R$ is $\Q$ with the
holes filled.

\smallskip
\textbf{The gap.} That narrative quietly presumes a notion of ``hole,'' and a
hole is only visible relative to a way of measuring closeness. Rudin uses
the ordinary absolute value throughout and never says it is a choice.

\smallskip
\textbf{Where the gap bites.} Nowhere in this book --- which is exactly why
it is worth one page. A reader who never learns that the choice existed will
later meet $\Q_p$ and find it inexplicable, and will also misread the word
``complete'' in Chapter 3 as a property of a \emph{set} rather than of a
\emph{metric}.

\smallskip
\textbf{What this section supplies.} A second absolute value on $\Q$, and the
different field it completes to.
\end{roadmap}

\sitem{S1.12}{Definition}
\emph{(The $p$-adic absolute value.) Fix a prime $p$. Every nonzero rational
$x$ is uniquely $x = p^{v}\,m/n$ with $m$, $n$ coprime to $p$; put
$v(x) = v$ and}
\[ |x|_p = p^{-v(x)}, \qquad d_p(x,y) = |x-y|_p . \]
\emph{Then $d_p$ is a metric on $\Q$, and $(\Q, d_p)$ is not complete; its
completion is the field $\Q_p$ of $p$-adic numbers.}
\provenance{Yu Pin p.\,125, as a worked corollary of his completion theorem
(Thm 77, pp.\,123--124).}

\begin{deeper}[What ``complete'' is a property of]
Under $d_p$ a number is \emph{small} when it is divisible by a high power of
$p$. With $p = 5$, the numbers $5, 25, 125, \dots$ march to zero, and
$1/5$ is enormous. The metric is genuinely a metric --- and it satisfies
something stronger than the triangle inequality,
$|x+y|_p \le \max(|x|_p,|y|_p)$, which makes its geometry unrecognisable:
every triangle is isosceles, and every point of a ball is its centre.

Complete $(\Q,d_p)$ and you get $\Q_p$, not $\R$. So:

\smallskip
\emph{``$\R$ is $\Q$ with the holes filled'' is incomplete as a description.}
The correct statement is that $\R$ is $\Q$ completed \emph{with respect to
the ordinary absolute value}. Change the absolute value and the holes move.

\smallskip
\emph{This is why Chapter 3 defines completeness for a metric space and not
for a set.} Rudin's 3.11 is stated for metric spaces precisely because the
property depends on $d$; his Exercise 3.24 constructs the completion in
general. The Chapter 2 appendix (S2.8) gives the definition, and Yu Pin's
universal property is what makes ``\emph{the} completion'' legitimate
language.

\smallskip
There is a theorem lurking here that neither book states: Ostrowski's, that
the ordinary absolute value and the $p$-adic ones are the \emph{only}
absolute values on $\Q$ up to equivalence. So $\R$ and the $\Q_p$ are the
complete list of completions --- Rudin's choice was not arbitrary, merely
unremarked.
\end{deeper}

% ============================================================
\ssection{Supplement G --- Real exponents, worked}

\begin{roadmap}[Why this supplement belongs to Chapter 1]
\raggedright
\textbf{What Rudin does.} Theorem 1.21 proves that every positive real has a
unique positive $n$th root. Exercises 1.6 and 1.7 then ask the reader to
build $b^x$ for all real $x$ and establish its laws, with a hint.

\smallskip
\textbf{The gap.} The exercises are substantially harder than they look, and
two steps in particular are glossed by the hint: that $b^{m/n}$ does not
depend on how the fraction $m/n$ is written, and that the limit defining
$b^x$ for irrational $x$ exists at all.

\smallskip
\textbf{Where the gap bites.} Every later appearance of $x^\alpha$,
$a^x$ and $\log$ --- 3.20's special limits, 8.6's exponential --- rests on
this exercise having been done. A reader who skipped it has a gap under the
whole of Chapters 3 and 8.

\smallskip
\textbf{What this section supplies.} Tao's worked version of Rudin's
exercise, with both glossed steps proved.
\end{roadmap}

\sitem{S1.13}{Theorem}
\emph{(Rational exponents.) For $x>0$ and $n$ a positive integer put
$x^{1/n} = \sup\{y \ge 0 : y^n \le x\}$. The set is nonempty and bounded
above, so the supremum exists; it satisfies $(x^{1/n})^n = x$, and for
integers $a$ and positive integers $b$ the definition
$x^{a/b} = (x^{1/b})^a$ is independent of the representation of the
fraction. With these, $x^{q+r} = x^qx^r$ and $(x^q)^r = x^{qr}$ for all
rational $q,r$.}
\provenance{Tao I, Def.\ 5.6.4--5.6.7 and Lemmas 5.6.5--5.6.9
(bk.\ pp.\,123--125) --- Rudin's Exercise 1.6 worked out.}
\begin{proof}
The set $\{y \ge 0 : y^n \le x\}$ contains $0$ and is bounded above by
$\max\{1,x\}$, since $y > \max\{1,x\}$ forces $y^n \ge y > x$. So the
supremum exists. Because $t \mapsto t^n$ is strictly increasing on
$[0,\infty)$, the set is the interval $[0,y]$ where $y$ is the unique
positive $n$th root supplied by 1.21; hence $x^{1/n}=y$ and
$(x^{1/n})^n = x$.

For independence, suppose $a/b = c/d$, so $ad = cb$. Both $(x^{1/b})^a$ and
$(x^{1/d})^c$ are positive, and raising each to the power $bd$ gives
$x^{ad}$ and $x^{cb}$ respectively --- equal numbers. Since 1.21 provides
\emph{exactly one} positive $bd$th root, the two agree.

The laws $x^{q+r}=x^qx^r$ and $(x^q)^r = x^{qr}$ now follow by writing $q$
and $r$ over a common denominator and using the corresponding integer
identities together with uniqueness in 1.21.
\end{proof}

\noindent
Rudin's 1.21 already gives the root; the content here is the
\emph{independence}, which needs the uniqueness half of 1.21 to force
$(x^{1/bd})^{ad} = (x^{1/b})^a$ whenever $a/b = c/d$.\kn{Rudin's hint for
Exercise 1.6 says to prove $(b^m)^{1/n} = (b^p)^{1/q}$ when $m/n = p/q$ and
then define $b^{m/n}$ to be that common value. That single sentence is the
independence step, and it is where all the work is --- both sides are $n$th
roots, so uniqueness in 1.21 is what identifies them.}

\sitem{S1.14}{Theorem}
\emph{(Irrational exponents.) Let $x>0$ and let $\alpha$ be real. If
$(q_n)$ is any sequence of rationals with $q_n \to \alpha$, then
$(x^{q_n})$ converges, and the limit does not depend on the sequence chosen.
One may therefore define $x^{\alpha} = \lim_n x^{q_n}$.}
\provenance{Tao I, Lemma 6.7.1 (bk.\ pp.\,152--153) --- Rudin's Exercise 1.7.}
\begin{proof}
Take $x \ge 1$ (for $0<x<1$ apply the result to $1/x$). If $q,r$ are
rationals with $|q|,|r| \le M$ and $|q-r| \le 1/K$, then
\[ \bigl|x^{q}-x^{r}\bigr| = x^{r}\bigl|x^{q-r}-1\bigr|
   \le x^{M}\bigl(x^{1/K}-1\bigr), \]
and $x^{1/K} \to 1$ as $K \to \infty$ by 3.20(b). Now let $q_n \to \alpha$
with each $q_n$ rational. The sequence $(q_n)$ is Cauchy and bounded, so the
displayed estimate makes $\bigl(x^{q_n}\bigr)$ Cauchy, and it converges by
completeness of $\R$ (3.11).

If $(q_n')$ is a second rational sequence tending to $\alpha$, interleave the
two into a single sequence; it also tends to $\alpha$, so its image
converges, and the two original limits --- being limits of subsequences of a
convergent sequence --- coincide.
\end{proof}

\begin{unpack}[The estimate that makes the definition legal]
Everything rests on one inequality. For rationals $q, r$ with
$|q - r| \le 1/K$ and both at most $M$ in absolute value,
\[ \bigl|x^{q} - x^{r}\bigr| \;\le\; x^{M}\bigl(x^{1/K} - 1\bigr)
   \qquad (x \ge 1), \]
and $x^{1/K} \to 1$ as $K \to \infty$ --- which is Rudin's own 3.20(b) in
disguise. So $(x^{q_n})$ is Cauchy whenever $(q_n)$ is, and completeness of
$\R$ supplies the limit. Interleaving two sequences with the same limit shows
the value is independent of the choice, exactly as in the uniqueness argument
of S1.1.

Two things are worth noticing. First, the definition needs completeness of
$\R$, so real exponentiation is downstream of 1.19 --- it could not have been
introduced in $\Q$. Second, this is the same three-step pattern used
throughout: \emph{approximate by rationals, show the approximations are
Cauchy, pass to the limit}. S1.1 uses it to build an isomorphism, S1.7 to
build a supremum, and Chapter 4's S4.17 to extend a uniformly continuous
function. Recognising the pattern is worth more than any one instance.
\end{unpack}
